% Exported from Math++. Compile with pdflatex.
\documentclass[11pt,a4paper]{article}
\usepackage[T1]{fontenc}
\usepackage[utf8]{inputenc}
\usepackage{lmodern}
\usepackage{amsmath,amssymb,amsthm,mathtools}
\usepackage{graphicx,booktabs,xcolor}
\usepackage[hidelinks]{hyperref}
\usepackage[margin=25mm]{geometry}
\theoremstyle{plain}
\newtheorem{theorem}{Theorem}
\newtheorem{lemma}[theorem]{Lemma}
\newtheorem{corollary}[theorem]{Corollary}
\newtheorem{proposition}[theorem]{Proposition}
\newtheorem{conjecture}[theorem]{Conjecture}
\theoremstyle{definition}
\newtheorem{definition}[theorem]{Definition}
\newtheorem{example}[theorem]{Example}
\theoremstyle{remark}
\newtheorem*{remark}{Remark}
\setlength{\parindent}{0pt}\setlength{\parskip}{0.6em}
\title{Piecewise functions, intervals and sets}
\date{}
\begin{document}
\maketitle
A function made of two pieces that meet smoothly at $x = 1$:

\begin{align*}
  f(x) \coloneqq  \begin{cases} x^2 & x < 1 \\ 2x - 1 & x \ge 1 \end{cases}
\end{align*}
\begin{equation*}
  f(3) = 5
\end{equation*}
The derivative works piece by piece, and the working notes what happens where the pieces meet:

\begin{align*}
  \frac{d}{dx} f(x) = \begin{cases} 2 x & x < 1 \\ 2 & x \ge 1 \end{cases}
\end{align*}
\begin{equation*}
  \int_0^2 f(x) \, dx = \frac{7}{3}
\end{equation*}
\begin{align*}
  \int f(x) \, dx = \begin{cases} \frac{x^{3}}{3} & x < 1 \\ x^{2} - x + \frac{1}{3} & x \ge 1 \end{cases} + C
\end{align*}
A step has a jump, so it is not differentiable there:

\begin{align*}
  h(x) \coloneqq  \begin{cases} 1 & x < 2 \\ 3 & x \ge 2 \end{cases}
\end{align*}
\begin{equation*}
  \int_0^4 h(x) \, dx = 8
\end{equation*}
\section*{Intervals and sets}

\begin{equation*}
  [0, 1) \cup (2, 3] = \left[0, 1\right) \cup \left(2, 3\right]
\end{equation*}
\begin{equation*}
  [0, 3] \cap (1, 5) = \left(1, 3\right]
\end{equation*}
\begin{equation*}
  \mathbb{R} \setminus [0, 1] = \left(-\infty, 0\right) \cup \left(1, \infty\right)
\end{equation*}
\begin{equation*}
  2 \in [0, 1) \cup (1, 3] = \text{true}
\end{equation*}
\begin{equation*}
  \{x \in \mathbb{R} : x^2 < 4\} = \left(-2, 2\right)
\end{equation*}
\begin{equation*}
  \{n \in \mathbb{Z} : n^2 < 10\} = \left\{-3, -2, -1, 0, 1, 2, 3\right\}
\end{equation*}
\section*{Inequalities}

Solve $x^2 - 4 < 0$ for $x$.

Solve $\frac{x - 1}{x + 2} \le 0$ for $x$.

Solve $x^2 \ge 4$ for $x$.

\end{document}
